%% ════════════════════════════════════════════════════════════════════════════════
%% Academic Poster — Tree of Thoughts for Automated Code Debugging
%% CS 4782, Cornell University · 36" × 24" landscape
%%
%% Compile: pdflatex poster.tex  (Overleaf: pdfLaTeX engine)
%% Upload to Overleaf: fix_rates.png, bug_type_breakdown.png
%%   (from results_new_run_charts/)
%% ════════════════════════════════════════════════════════════════════════════════
\documentclass[final]{beamer}

\usepackage[T1]{fontenc}
\usepackage[utf8]{inputenc}
\usepackage[
  orientation=landscape,
  size=custom,
  width=91.44,
  height=60.96,
  scale=1.2
]{beamerposter}

\usepackage{lmodern}
\usepackage{booktabs}
\usepackage{amsmath}
\usepackage{graphicx}
\graphicspath{{../results_new_run_charts/}}
\usepackage{tikz}
\usepackage{xcolor}

%% ── Colours ───────────────────────────────────────────────────────────────────
\definecolor{cornellred}{RGB}{179,27,27}
\definecolor{lightgray}{RGB}{245,245,245}
\definecolor{accentblue}{RGB}{30,90,180}
\definecolor{accentgreen}{RGB}{40,150,80}
\definecolor{accentpurple}{RGB}{120,40,180}

%% ── Beamer theme ──────────────────────────────────────────────────────────────
\usetheme{default}
\setbeamercolor{background canvas}{bg=white}

\setbeamercolor{block title}{fg=white, bg=accentblue}
\setbeamercolor{block body}{fg=black, bg=lightgray}
\setbeamercolor{block title alerted}{fg=white, bg=cornellred}
\setbeamercolor{block body alerted}{fg=black, bg=cornellred!10}
\setbeamercolor{block title example}{fg=white, bg=accentgreen}
\setbeamercolor{block body example}{fg=black, bg=accentgreen!10}

\setbeamerfont{block title}{size=\large, series=\bfseries}
\setbeamerfont{block body}{size=\normalsize}

%% ════════════════════════════════════════════════════════════════════════════════
\begin{document}
\begin{frame}[t]{}

%% ── HEADER (reduced to ~6 cm to leave more content space) ────────────────────
\begin{tikzpicture}[remember picture, overlay]
  \fill[cornellred] (current page.north west)
    rectangle ([yshift=-6.5cm]current page.north east);
\end{tikzpicture}

\vspace{0.3cm}
\centering
{\fontsize{42}{50}\bfseries\color{white}\selectfont
  Tree of Thoughts for Automated Code Debugging\par}
\vspace{0.4cm}
{\fontsize{22}{28}\color{white}\selectfont
  Ronald Feng \quad Jefferson Zhou \quad Haochen Wang
  \quad\textbullet\quad Cornell University
  \quad\textbullet\quad CS\,4782: Intro to Deep Learning\par}
\vspace{0.5cm}

%% ── THREE-COLUMN BODY ────────────────────────────────────────────────────────
\begin{columns}[t,totalwidth=\textwidth]

%% ══════════════════════ COLUMN 1 — Intro + Method ════════════════════════════
\begin{column}{0.29\textwidth}

\begin{block}{Introduction \& Motivation}
\textbf{Problem:} LLMs struggle with debugging — it requires exploring \emph{multiple} hypotheses and backtracking from wrong diagnoses. Standard IO / CoT locks into one path.

\vspace{0.4cm}
\textbf{Solution:} \textbf{Tree of Thoughts (ToT)} (Yao et al., NeurIPS 2023) lets an LLM search a tree of reasoning steps. We re-implement and extend it with \textbf{MCTS} for code debugging.

\vspace{0.4cm}
\textbf{Key advantage of code debugging as a domain:}
\begin{itemize}
\setlength{\itemsep}{2pt}
  \item Test pass/fail provides \textbf{unambiguous reward} — no noisy LLM self-evaluation needed
  \item Bugs are localized to single functions — tractable search space
  \item Automatic, reproducible evaluation
\end{itemize}

\vspace{0.4cm}
\textbf{Target result:} Table 1 of Yao et al. — ToT-BFS achieves \textbf{74\%} on Game of 24 vs \textbf{4\%} for IO/CoT (18$\times$ gain). We reproduce this structured-search advantage in debugging.
\end{block}

\vspace{0.5cm}

\begin{block}{Dataset \& Setup}
\textbf{HumanEval-Bugs} — 164 Python problems

\vspace{0.35cm}
\begin{tabular}{@{} l c @{}}
\toprule
\textbf{Bug Category} & \textbf{Count} \\
\midrule
\texttt{wrong\_operator}    & 106 \\
\texttt{incorrect\_return}  & 33  \\
\texttt{missing\_condition} & 17  \\
\texttt{off\_by\_one}       & 8   \\
\midrule
Total                        & 164 \\
\bottomrule
\end{tabular}

\vspace{0.35cm}
\textbf{LLM:} Qwen2.5-Coder-7B-Instruct (Ollama)\\
\textbf{Cost:} \$0 — open-source, local inference\\
\textbf{Hardware:} Google Colab (NVIDIA GPU)\\
\textbf{Eval:} Sandboxed \texttt{subprocess}, 5\,s timeout
\end{block}

\vspace{0.5cm}

\begin{block}{Methodology}
\textbf{Shared 4-component pipeline:}
\begin{enumerate}
\setlength{\itemsep}{2pt}
  \item \textbf{Thought Generator} — LLM produces $k{=}3$ bug-hypothesis candidates
  \item \textbf{Thought Evaluator} — Hybrid: 60\% LLM score + 40\% execution signal
  \item \textbf{Search Controller} — BFS, DFS, or MCTS
  \item \textbf{Fix Generator} — Per hypothesis: generate $k$ fixes, verify by tests
\end{enumerate}

\vspace{0.35cm}
\textbf{Modifications vs.\ original ToT:}
\begin{itemize}
\setlength{\itemsep}{2pt}
  \item Execution-based reward replaces LLM self-evaluation
  \item Two-level tree: hypotheses $\to$ fixes
  \item \textbf{MCTS} added as novel third search strategy
\end{itemize}
\end{block}

\end{column}

%% ══════════════════════ COLUMN 2 — Results ═══════════════════════════════════
\begin{column}{0.42\textwidth}

\begin{exampleblock}{Results: Fix Rate by Method (n = 164)}
\centering
\includegraphics[width=0.97\linewidth]{fix_rates.png}

\vspace{0.3cm}
\raggedright
\begin{tabular}{@{} l c c c c @{}}
\toprule
\textbf{Method} & \textbf{Fix Rate} & \textbf{1st-Attempt} & \textbf{Avg Tokens} & \textbf{Avg Time} \\
\midrule
IO                        & 87.8\%           & 87.8\%               & 567                 & 2.0\,s \\
CoT                       & 77.4\%           & 77.4\%               & 940                 & 5.0\,s \\
CoT-SC ($n{=}5$)          & 95.7\%           & 95.7\%               & 4,749               & 24.6\,s \\
ToT-BFS ($k{=}3$)         & 98.8\%           & 79.9\%               & 2,874               & 14.7\,s \\
ToT-DFS ($k{=}3$)         & 98.8\%           & 86.6\%               & 2,798               & 14.2\,s \\
\textbf{MCTS} ($k{=}3$, $s{=}12$) & \textbf{100\%} & \textbf{92.1\%} & 14,335 & 64.0\,s \\
\bottomrule
\end{tabular}
\end{exampleblock}

\vspace{0.5cm}

\begin{alertblock}{Novel Contribution: MCTS Search}
\textbf{Why MCTS?} BFS exhausts all nodes; DFS commits early. MCTS \emph{learns} which hypotheses deserve more exploration via UCB1.

\vspace{0.35cm}
\textbf{UCB1 node selection:}
\[
  \mathrm{UCB1}(v) = \underbrace{\frac{w_v}{n_v}}_{\text{exploit}}
  + C\sqrt{\frac{\ln n_{\mathrm{parent}}}{n_v}}, \quad C = \sqrt{2}
\]

\textbf{Four phases per iteration:}
\begin{enumerate}
\setlength{\itemsep}{1pt}
  \item \textbf{Select} via UCB1 to most promising leaf
  \item \textbf{Expand} — new hypothesis if leaf already visited
  \item \textbf{Simulate} — generate fix, run tests $\to$ reward $\in \{0, 0.1, 1.0\}$
  \item \textbf{Backpropagate} — update $w_v$, $n_v$ to root
\end{enumerate}

\vspace{0.3cm}
\textcolor{accentpurple}{\textbf{Result: MCTS is the only method to achieve 100\% fix rate.}} With $s{=}12$ simulations it explores 46.8 nodes on average, allocating more rollouts to promising hypotheses rather than uniformly expanding all branches.
\end{alertblock}

\end{column}

%% ══════════════════════ COLUMN 3 — Analysis + Insights ══════════════════════
\begin{column}{0.26\textwidth}

\begin{block}{Fix Rate by Bug Type (ToT-BFS)}
\centering
\includegraphics[width=0.97\linewidth]{bug_type_breakdown.png}

\raggedright\vspace{0.2cm}
\texttt{off\_by\_one} (87.5\%) is hardest — ambiguous boundary requires multiple fix attempts.
All other categories hit $\geq$99\%.
\end{block}

\vspace{0.5cm}

\begin{block}{Key Insights}
\begin{itemize}
\setlength{\itemsep}{4pt}
  \item \textbf{CoT hurts vs.\ IO} (77.4\% vs.\ 87.8\%) — Qwen2.5-Coder-7B \emph{over-reasons}: extended chain-of-thought introduces errors that direct prompting avoids.

  \item \textbf{ToT search recovers} — BFS/DFS reach 98.8\% by generating multiple fix candidates and verifying each against tests.

  \item \textbf{MCTS achieves 100\%} — execution-based backpropagation corrects noisy initial LLM scores; poor hypotheses are deprioritised after failed rollouts.

  \item \textbf{Token tradeoff} — MCTS uses $5\times$ more tokens than BFS/DFS. Since Qwen runs locally at \$0, cost is wall-clock time (${\sim}$64\,s/problem), not dollars.

  \item \textbf{Backtracking matters} — DFS with backtrack fixes 84.6\% of backtracked problems vs.\ 100\% without, confirming that initial hypotheses are often wrong.
\end{itemize}
\end{block}

\vspace{0.5cm}

\begin{block}{Conclusion \& Future Work}
\textbf{ToT + MCTS achieves perfect 100\% fix rate} on HumanEval-Bugs using a free, open-source 7B model. Execution-based reward backpropagation is the decisive advantage over all baselines.

\vspace{0.35cm}
\textbf{Future directions:}
\begin{itemize}
\setlength{\itemsep}{2pt}
  \item Scale to repository-level bugs (multi-file context)
  \item Dynamic analysis: inject stack traces / runtime values
  \item Evaluate on DebugBench (4,253 bugs, 3 languages)
  \item Reduce MCTS token cost via early-stopping rollouts
\end{itemize}

\vspace{0.35cm}
\textbf{References:}\,\small
[1] Yao et al., \textit{Tree of Thoughts}, NeurIPS 2023.
[2] Chen et al., \textit{HumanEval}, arXiv 2021.
[3] Wang et al., \textit{Self-Consistency}, ICLR 2023.
[4] Coulom, \textit{MCTS in Games}, CG 2006.
\end{block}

\end{column}
\end{columns}

\end{frame}
\end{document}
